机器学习最容易从算法名称开始：先问该用线性模型、树还是神经网络，再把数据交给优化器。这个顺序把真正决定结论含义的那几步藏了起来。模型训练之前得先说清预测发生在谁身上、那一刻能看见什么、输出会触发什么行动、错误由谁承担。这些都写清楚之后，才轮到算法。

本单元要立住的判断有四条，后面所有章节都要在它们之下工作。

第一条，学习问题的第一步是把业务目标翻译成可优化的目标，而这一步最容易出错。凸优化里那句``目标写歪了，解会精确地钻那个空子''在这里同样成立，还多一层麻烦：优化器会一边钻空子，一边给你一条漂亮的下降曲线。指标选错、口径不一致、离线目标与线上目标不对齐，这些都不是模型问题，换模型也修不好。

第二条，泛化是关于未见数据的断言，所以评估必须真的模拟``未见''这件事。任何在划分之前动过全量数据的步骤——标准化、特征筛选、降维、重采样——都会破坏这个模拟，而且不会报错。

第三条，单一指标掩盖分布。准确率在类别不平衡下几乎不携带信息；精确率与召回率的取舍由代价比决定，不由模型决定；分组指标甚至可能与总体指标方向相反。只报一个数字的评估表，读者无法排除这些情形。

第四条，没有朴素基线的比较不构成证据。常数基线、随机基线、持续性基线、单特征基线、当前线上规则，都必须走完全相同的划分与预处理，差距还要放进置信区间里看。

\subsection{怎么读这一章}

第一次读，抓住``单位—时点—行动—代价''这条链，并且能区分损失、风险、训练目标和报告指标这四个量。只想尽快建立评价纪律的话，第一、二、四篇是最短路线；第三篇在需要追查``为什么回归用均方误差、分类用交叉熵''时再细读。读到能把一个含糊的``准确率问题''改写成一句带条件的陈述——在哪个采样机制下、对哪个预测单位、用什么损失、经过多少次选择、估计了哪个分布上的表现——就可以往后走。

\subsection{本章篇目}

以下四篇按概念依赖排列；若从具体困惑进入，也可以先打开最贴近的一篇，再沿篇内链接回补。

\begin{itemize}
\tightlist
\item \hyperref[sec:13-Machine-Learning/01-Learning-Problems-and-Evaluation/01]{``从现实决策到学习问题''}：固定预测单位、可用信息、输出与行动，再定义数据、假设空间与部署分布；
\item \hyperref[sec:13-Machine-Learning/01-Learning-Problems-and-Evaluation/02]{``损失、风险、目标与指标不是同一个量''}：四个量在代码里只隔几行，在数学上分住四层，中间隔着三次可能出错的翻译；
\item \hyperref[sec:13-Machine-Learning/01-Learning-Problems-and-Evaluation/03]{``目标函数与噪声模型的对应：从生成机制到优化目标''}：损失是把噪声密度取负对数翻过来写，形状决定优化难度，统计性质决定泛化；
\item \hyperref[sec:13-Machine-Learning/01-Learning-Problems-and-Evaluation/04]{``数据划分、泄漏与分布偏移''}：划分策略、泄漏的三种长法、交叉验证估计的对象、置信区间，以及什么时候必须做在线实验。
\end{itemize}

这一单元不比较任何具体模型，它只规定后面所有比较必须遵守的规矩。
